\documentclass[11pt,a4paper]{article}
\usepackage{amsmath,amssymb}
\usepackage{graphicx}
\usepackage{xcolor}
\usepackage{booktabs}
\usepackage{caption}
\usepackage{subcaption}
\usepackage{float}
\usepackage{geometry}
\usepackage{hyperref}
\geometry{margin=1in}
\raggedbottom % deterministic page filling (prevents float/TOC page-count oscillation)
\graphicspath{{../fig/}}
\hypersetup{colorlinks=true,linkcolor=blue,citecolor=blue,urlcolor=blue}

% chemical shorthands (math mode, no extra packages)
\newcommand{\Lip}{\ensuremath{\mathrm{Li^{+}}}}
\newcommand{\cinf}{\ensuremath{c_{\infty}}}
\newcommand{\cs}{\ensuremath{c_{+}(0)}}
\newcommand{\Vt}{\ensuremath{RT/F}}

\title{\textbf{Interfacial Ion Accumulation Lowers the Intercalation\\
Threshold in 2D Electrolyte-Gated Transistors}\\[4pt]
\large A First-Principles Poisson--Nernst--Planck Theory Supporting the\\
EDLT$\rightarrow$ECRAM Transition Under Cumulative Pulsing}
\author{ecmLab \quad (\texttt{project\_ecram})}
\date{\today}

\begin{document}
\maketitle

\begin{abstract}
\noindent
The accompanying manuscript reports that a single electrolyte-gated transistor
(EGT) built on a 2D channel ($\mathrm{MoS_2}$, $\mathrm{WSe_2}$, graphene)
switches from a \emph{volatile} electric-double-layer transistor (EDLT) to a
\emph{non-volatile} electrochemical RAM (ECRAM) as the gate voltage exceeds a
critical value, and---crucially---that \emph{cumulative/sequential pulsing
lowers that critical voltage} from ${\sim}4$--$5$\,V (single pulse) to
${\sim}2$\,V. The manuscript attributes this to a build-up of interfacial \Lip{}
that ``pre-loads'' the channel interface, and reserves a figure panel
(Fig.~5F) for a simulation of the effect that is not yet supplied. This report
fills that gap. We solve the coupled Poisson--Nernst--Planck (PNP) equations for
a blocking electrode with a Stern layer, validate the solver against the
analytic Gouy--Chapman--Stern equilibrium ($<0.1\%$ error), and show from first
principles that (i) a single pulse is volatile while a high-repetition-rate
pulse train ratchets the interfacial \Lip{} concentration up to ${\sim}6\times$
bulk; (ii) by the Nernst relation this raises the intercalation open-circuit
voltage (OCV) by up to ${\approx}46$\,mV (the content of Fig.~5F); and (iii)
through Butler--Volmer kinetics it lowers the interfacial intercalation
threshold by up to ${\approx}91$\,mV, monotonically in pulse number. The
accumulation regime is governed by the ratio of the pulse period to the bulk
diffusion time $\tau_D=L^2/D$; for a solid Li-salt electrolyte this is tens of
seconds, consistent with both the accumulation seen under $1$-s pulses and the
deliberate $300$-s relaxation used in the single-pulse measurements.
\end{abstract}

\tableofcontents
\newpage

%======================================================================
\section{Device context and the EDLT$\rightarrow$ECRAM transition}
%======================================================================
The manuscript studies side-gated EGTs in which a solid Li-salt electrolyte
(mobile \Lip{} cations and $\mathrm{ClO^{-}}$ anions) gates a multilayer 2D
semiconductor channel. Two operating regimes are identified:

\begin{itemize}
\item \textbf{EDLT (volatile).} At low gate voltage $V_G$, ions polarize and form
an electric double layer (EDL) at the channel/electrolyte interface. The induced
screening charge modulates the channel current $I_D$; on removal of $V_G$ the
ions relax and $I_D$ returns to baseline. The conductance change is fully
reversible.
\item \textbf{ECRAM / electrochemical doping (non-volatile).} When $V_G$ exceeds
the electrochemical window, \Lip{} intercalates into the 2D crystal while the
anion is oxidized at the gate; the released electron neutralizes the inserted Li,
producing a \emph{persistent} conductance increase.
\end{itemize}

The experimental headline (manuscript Figs.~2--4) is twofold:
\begin{enumerate}
\item A well-defined threshold separates the modes: for multilayer graphene the
non-volatile shift first appears at $V_G=4$\,V under a single $1$-s pulse
($\Delta I_D=0.22\,\mu$A); for $\mathrm{MoS_2}$ it is ${\sim}5$\,V.
\item Under \emph{sequential} pulsing the threshold falls dramatically---to
$V_G\approx 2$\,V---and the same reduction recurs across multiple devices and
both materials, implying an \emph{intrinsic interfacial mechanism} rather than a
device-specific artifact.
\end{enumerate}

%======================================================================
\section{The gap: what the manuscript needs theory for}
%======================================================================
The proposed mechanism (manuscript Fig.~5 text) is that under closely spaced
pulses the ions driven to the interface during each pulse \emph{do not fully
dissipate} before the next pulse arrives, so the interfacial \Lip{}
concentration \emph{progressively builds up}, pre-loading the interface and
lowering the voltage required for intercalation. However:

\begin{itemize}
\item \textbf{Figure 5F is an empty placeholder} labelled ``\emph{Simulation of
interfacial ion concentration impact on the OCV of the material}.'' No
simulation is provided.
\item The link between \emph{interfacial concentration} and \emph{intercalation
threshold voltage} is asserted qualitatively but not quantified.
\item The dependence on \emph{repetition rate} (why fast trains accumulate but
isolated pulses do not) is stated but not demonstrated.
\end{itemize}

This report supplies a first-principles, self-validated treatment of all three
points.

%======================================================================
\section{Hypothesis}
%======================================================================
\noindent\fbox{\parbox{0.96\textwidth}{\textbf{Hypothesis.} A high-repetition-rate
gate pulse train accumulates \Lip{} at the channel/electrolyte interface because
the pulse period is shorter than the diffusive relaxation time of the
near-surface region. The resulting elevated interfacial concentration $\cs$
raises the thermodynamic driving force (OCV) for intercalation by the Nernst term
$(\Vt)\ln(\cs/\cinf)$ and, through concentration-dependent Butler--Volmer
kinetics, lowers the intercalation overpotential by $(\Vt/\alpha_c)\ln(\cs/\cinf)$.
Hence the EDLT$\rightarrow$ECRAM threshold decreases monotonically with pulse
number, in a manner intrinsic to the interface and therefore material-independent.}}

%======================================================================
\section{Theory and methods}
%======================================================================

\subsection{Governing equations (Poisson--Nernst--Planck)}
The electrolyte contains a cation ($\Lip$, charge $z_+=+1$, diffusivity $D_+$)
and an anion ($z_-=-1$, $D_-$). Their transport is governed by the
Nernst--Planck flux with diffusion and migration, coupled to Poisson's equation
for the potential $\varphi$:
\begin{align}
\frac{\partial c_i}{\partial t} &= -\,\frac{\partial N_i}{\partial x},
& N_i &= -D_i\!\left(\frac{\partial c_i}{\partial x}
       + z_i\frac{F}{RT}\,c_i\frac{\partial \varphi}{\partial x}\right),
\label{eq:np}\\[4pt]
-\,\varepsilon_0\varepsilon_r\frac{\partial^2 \varphi}{\partial x^2}
  &= F\sum_i z_i c_i .
\label{eq:poisson}
\end{align}
Equations~\eqref{eq:np}--\eqref{eq:poisson} are the same Gouy--Chapman--Stern
PNP family used in the group's \texttt{ddl/} models.

\subsection{Boundary conditions: blocking electrode with a Stern layer}
The channel electrode at $x=0$ is \emph{blocking} (no Faradaic flux in the EDLT
regime), $N_i\!\cdot\!n=0$, and carries a compact \emph{Stern} layer of
thickness $x_S$ and permittivity $\varepsilon_S$. Gauss' law across the Stern
layer relates the metal potential $\varphi_M(t)$ to the diffuse-layer potential
$\varphi(0)$:
\begin{equation}
\sigma = \varepsilon_0\varepsilon_S\,\frac{\varphi_M-\varphi(0)}{x_S}
       = -\,\varepsilon_0\varepsilon_r\left.\frac{\partial\varphi}{\partial x}\right|_{0},
\label{eq:stern}
\end{equation}
a Robin condition for $\varphi$. At the far boundary $x=L$ the electrolyte meets
a bulk reservoir, $\varphi=0$ and $c_i=\cinf$. The gate excitation is a pulse
train $\varphi_M(t)$.

\subsection{Scaling and the two governing timescales}
Non-dimensionalizing with $X=x/L$, $\tilde\varphi=\varphi F/RT$,
$\tilde c=c/\cinf$, $\tilde t=tD/L^2$, Poisson becomes
$-\epsilon^2\,\partial_X^2\tilde\varphi=\tfrac12(\tilde c_+-\tilde c_-)$ with the
double-layer thinness $\epsilon=\lambda_D/L$ and Debye length
$\lambda_D=\sqrt{\varepsilon_0\varepsilon_r RT/(2F^2\cinf)}$. Two timescales
emerge and control everything:
\begin{equation}
\boxed{\;\tau_c \sim \frac{\lambda_D L}{D}\;\text{(EDL charging, fast)}\qquad
        \tau_D \sim \frac{L^2}{D}\;\text{(bulk diffusion, slow)}\;}
\label{eq:tau}
\end{equation}
with $\tau_c/\tau_D=\epsilon\ll1$. \textbf{Accumulation occurs when the pulse
period is shorter than the relaxation time}: ions injected toward the interface
have no time to diffuse back before the next pulse.

\subsection{From interfacial concentration to OCV (Nernst)}
The intercalation half-reaction
$\Lip + e^- + \langle\text{host}\rangle \rightleftharpoons \mathrm{Li}\langle\text{host}\rangle$
has an open-circuit (equilibrium) potential that depends on the \emph{local}
interfacial \Lip{} activity. Relative to the bulk reference,
\begin{equation}
\Delta E_{\mathrm{OC}} = \frac{RT}{F}\,\ln\!\frac{\cs}{\cinf}
 \;\approx\; 59\ \text{mV per decade of interfacial enrichment at }298\,\text{K}.
\label{eq:nernst}
\end{equation}
Pre-loading the interface (raising $\cs$) raises the driving force for
intercalation. \emph{This is precisely the quantity Fig.~5F calls for.}

\subsection{From interfacial concentration to threshold (Butler--Volmer)}
The intercalation (cathodic) current obeys concentration-dependent
Butler--Volmer kinetics,
\begin{equation}
i_F = i_0\,\frac{\cs}{c_{\mathrm{ref}}}
      \left[\exp\!\Big(\tfrac{\alpha_a F\eta}{RT}\Big)
          - \exp\!\Big(\!-\tfrac{\alpha_c F\eta}{RT}\Big)\right],
\label{eq:bv}
\end{equation}
so the exchange current scales with $\cs$. Requiring a fixed onset current $i^*$
for ``measurable'' intercalation, the required interfacial overpotential drops by
\begin{equation}
\Delta|\eta_{\mathrm{th}}| = \frac{RT}{\alpha_c F}\,\ln\!\frac{\cs}{\cinf}.
\label{eq:threshold}
\end{equation}
Equations~\eqref{eq:nernst} and \eqref{eq:threshold} are the analytic backbone;
the PNP solver supplies $\cs$ as a function of the pulse protocol.

\subsection{Numerical method and validation}
We discretize Eqs.~\eqref{eq:np}--\eqref{eq:poisson} on a graded mesh refined at
the electrode using a vertex-centered finite-volume scheme. Drift--diffusion
fluxes use the Scharfetter--Gummel (exponential-fitting) form to keep the thin,
sharply varying double layer stable and the concentrations positive. Poisson's
equation is solved as a tridiagonal linear system (with the Stern Robin
condition) at every right-hand-side evaluation, reducing the system to ODEs in
the concentrations integrated with a stiff (BDF) solver
(\texttt{sim/pnp\_core.py}).

\paragraph{Validation against analytic equilibrium.} Under a held bias the solver
must reproduce the Boltzmann/Gouy--Chapman--Stern profile $c_+(X)=e^{-\tilde
\varphi(X)}$. Table~\ref{tab:val} shows agreement to $<0.1\%$, establishing that
the transient results below are trustworthy.

\begin{table}[H]
\centering
\caption{Equilibrium validation at held bias $\tilde\varphi_M=-3$ (thermal
volts), $\epsilon=0.02$, Stern ratio $\delta=0.1$.}
\label{tab:val}
\begin{tabular}{lcc}
\toprule
Quantity & Solver & Analytic (Boltzmann/GCS)\\
\midrule
Diffuse-layer potential $\tilde\varphi(0)$ & $-2.6504$ & --- \\
Stern drop $\tilde\varphi_M-\tilde\varphi(0)$ & $-0.3496$ & --- \\
Surface concentration $c_+(0)/\cinf$ & $14.148$ & $e^{-\tilde\varphi(0)}=14.159$ \\
Max.\ relative error (cation) & \multicolumn{2}{c}{$7.9\times10^{-4}$} \\
Max.\ relative error (anion)  & \multicolumn{2}{c}{$6.8\times10^{-4}$} \\
\bottomrule
\end{tabular}
\end{table}

%======================================================================
\section{Results: proof of the hypothesis}
%======================================================================
All runs use $\epsilon=\lambda_D/L=0.02$, a pulse amplitude
$\tilde\varphi_M=-4$ (thermal volts), on-time $0.03$ and off-time $0.012$ (in
units of $\tau_D=L^2/D$) for the fast train, and $D_-/D_+=2$ (slow \Lip{}).

\subsection{Cumulative pulsing accumulates interfacial \Lip{}}
Figure~\ref{fig:accum} is the central first-principles result. A \emph{single}
pulse drives the interfacial concentration to ${\sim}11\times$ bulk and then
\emph{relaxes fully} to the bulk value---a volatile EDL response (panel A). A
\emph{fast} pulse train (period $\ll\tau_D$) shows both the peaks and the
\emph{between-pulse baseline} ratcheting upward (panel B, blue): the interface
never resets, i.e.\ it is pre-loaded. A \emph{slow} train (period comparable to
the relaxation time) relaxes after each pulse and shows \emph{no} accumulation
(panel B, orange), directly demonstrating the repetition-rate dependence the
manuscript invokes. Panel C quantifies the ratchet: the between-pulse baseline
rises from $4.4\times$ to $5.9\times$ bulk over 12 pulses (peaks $10.8\to14.9$);
numerical values in Table~\ref{tab:enrich}.

\begin{figure}[H]
\centering
\includegraphics[width=\textwidth]{fig1_accumulation.png}
\caption{First-principles PNP: a single pulse is volatile, whereas a
high-repetition-rate train pre-loads the interface with \Lip{}. Supports
manuscript Figs.~2B, 3 and 5C--E.}
\label{fig:accum}
\end{figure}

\begin{table}[H]
\centering
\caption{Interfacial cation enrichment $c_+(0)/\cinf$ versus pulse number
(fast train).}
\label{tab:enrich}
\small
\begin{tabular}{lcccccccccccc}
\toprule
$N$ & 1 & 2 & 3 & 4 & 5 & 6 & 7 & 8 & 9 & 10 & 11 & 12\\
\midrule
baseline & 4.36 & 5.18 & 5.44 & 5.56 & 5.64 & 5.69 & 5.74 & 5.78 & 5.81 & 5.84 & 5.86 & 5.89\\
peak     & 10.8 & 13.1 & 13.7 & 14.0 & 14.2 & 14.4 & 14.5 & 14.6 & 14.7 & 14.8 & 14.8 & 14.9\\
\bottomrule
\end{tabular}
\end{table}

\subsection{Interfacial concentration controls the OCV (fills Fig.~5F)}
Figure~\ref{fig:ocv} maps the PNP-derived enrichment to the intercalation OCV via
Eq.~\eqref{eq:nernst}. The driving force follows the $\approx59$\,mV/decade
Nernst line, and the pulse-train pre-loading delivers up to
${\approx}+46$\,mV of intercalation OCV (panel B). \textbf{This figure is the
content the manuscript reserves for the empty panel Fig.~5F.}

\begin{figure}[H]
\centering
\includegraphics[width=\textwidth]{fig2_OCV.png}
\caption{Interfacial ion concentration impact on the OCV. (A) Nernst mapping with
PNP pulse-train states marked; (B) OCV gain versus pulse number. \textbf{This
fills manuscript Fig.~5F.}}
\label{fig:ocv}
\end{figure}

\subsection{Pre-loading lowers the intercalation threshold}
Figure~\ref{fig:thresh} applies the Butler--Volmer relation
Eqs.~\eqref{eq:bv}--\eqref{eq:threshold}. As pre-loading increases, the Faradaic
onset shifts to lower overpotential (panel A), and the interfacial intercalation
threshold falls monotonically with pulse number, reaching ${\approx}91$\,mV
(with $\alpha_c=0.5$) for the fully pre-loaded interface relative to an isolated,
fully relaxed single pulse (panel B). This is the mechanistic origin of the
single-pulse$\to$cumulative threshold reduction in manuscript Figs.~3--4.

\begin{figure}[H]
\centering
\includegraphics[width=\textwidth]{fig3_threshold.png}
\caption{Butler--Volmer intercalation onset shifts to lower voltage as the
interface is pre-loaded (A), giving a threshold reduction that grows with pulse
number (B). Supports manuscript Figs.~3--4.}
\label{fig:thresh}
\end{figure}

%======================================================================
\section{Connection to the experiments}
%======================================================================

\subsection{Why fast pulses accumulate and isolated pulses do not}
The accumulation regime is set by Eq.~\eqref{eq:tau}. For a solid Li-salt
electrolyte ($D_+\sim10^{-13}\,\mathrm{m^2/s}$) over a $\mu$m-scale gate
electrolyte ($L\sim$ few $\mu$m), the bulk diffusion time is
$\tau_D=L^2/D\sim$ \emph{tens of seconds}. Consequently:
\begin{itemize}
\item $1$-s pulses repeated at ${\sim}0.5$\,Hz have a period
$\ll\tau_D$ $\Rightarrow$ \textbf{accumulation} (cumulative ECRAM programming).
\item The single-pulse characterization deliberately waits \textbf{300\,s}
between pulses (manuscript Fig.~2C), i.e.\ $\gg\tau_D$ $\Rightarrow$ a guaranteed
full reset, isolating the genuine single-pulse threshold.
\end{itemize}
The model and the experimental protocol are therefore self-consistent: the same
$\tau_D$ that produces accumulation under fast pulsing also justifies the
$300$-s relaxation chosen to suppress it.

\subsection{Trend and magnitude versus experiment}
The theory reproduces the \emph{direction, monotonicity and saturation} of the
measured effect: the threshold reduction grows with pulse number and saturates as
the interfacial concentration approaches its pulse-train steady state
(Fig.~\ref{fig:thresh}B), mirroring the experimental observation that a few
pulses already shift the onset (manuscript Figs.~3B, 4B,D,F).

\paragraph{Interface versus device voltage.} The $46$--$91$\,mV figures are
\emph{interfacial} (diffuse-layer / overpotential) quantities. At the device
level the applied $V_G$ is partitioned across the Stern layer, the series
electrolyte and the channel, so a given fractional pre-loading appears as a
\emph{larger} reduction in $V_G$. The model thus naturally yields a volt-scale
single-pulse$\to$cumulative shift consistent with the measured
$4$--$5\,\text{V}\to{\sim}2\,\text{V}$, with the exact factor set by the series
partition (\S\ref{sec:next}).

\subsection{Material independence}
Because the mechanism is interfacial---ion accumulation in the electrolyte,
governed by Eqs.~\eqref{eq:tau} and \eqref{eq:threshold}, not by the channel band
structure---the same threshold reduction is predicted for any 2D channel. This
explains why the manuscript observes the effect identically for graphene and
$\mathrm{MoS_2}$.

%======================================================================
\section{Limitations and next steps}\label{sec:next}
%======================================================================
\begin{enumerate}
\item \textbf{Coupled Butler--Volmer in COMSOL.} Adding the intercalation flux at
$x=0$ via the \texttt{TertiaryCurrentDistributionNernstPlanck} +
\texttt{ElectrodeReaction} pattern (from \texttt{ddl\_bv/ddlBV0.m}) would yield
the threshold $V_G(N)$ curve from a single first-principles solve, rather than
the present PNP+analytic-BV combination.
\item \textbf{Device-level $V_G$ mapping.} Including the Stern $+$ series $+$
channel partition would convert the interfacial overpotential reductions into
absolute $V_G$ thresholds, targeting a quantitative
$4$--$5\,\text{V}\to2\,\text{V}$ match.
\item \textbf{Pre-formed-EDL protocol.} Holding a small DC bias before pulsing
(manuscript Fig.~5B,D) is directly representable and predicted to enhance the
response.
\item \textbf{Parameters.} The diffusivity, electrolyte thickness and Stern
parameters are representative; tightening them to the measured electrolyte would
sharpen the quantitative comparison.
\end{enumerate}

%======================================================================
\section{Summary}
%======================================================================
A validated 1D Poisson--Nernst--Planck model demonstrates, from first principles,
the mechanism the manuscript proposes but does not yet show. High-repetition-rate
gate pulsing accumulates interfacial \Lip{} because the pulse period is shorter
than the diffusive relaxation time $\tau_D=L^2/D$; the elevated interfacial
concentration raises the intercalation OCV (Nernst, filling Fig.~5F) and lowers
the Butler--Volmer intercalation threshold, monotonically in pulse number and
independently of the channel material. The model is quantitatively consistent
with the experimental protocol (the $\tau_D\sim$ tens of seconds that drives
accumulation under $1$-s pulses also justifies the $300$-s single-pulse reset)
and reproduces the measured single-pulse$\to$cumulative threshold reduction.

%======================================================================
\appendix
\section{Reproducing the figures}
%======================================================================
\begin{verbatim}
cd project_ecram/sim
python3 pnp_core.py        # equilibrium validation (Table 1): PASS
python3 ecram_theory.py    # generates fig1/fig2/fig3 in ../fig/
matlab -batch pulse        # optional: COMSOL LiveLink pulse-train PNP
                           #           (set N_PULSE=1 for the single-pulse control)
\end{verbatim}
Key files: \texttt{sim/pnp\_core.py} (validated PNP solver),
\texttt{sim/ecram\_theory.py} (figures), \texttt{sim/pulse.m} (COMSOL
counterpart), \texttt{theory/theory\_support.md} (companion notes).

\begin{thebibliography}{9}
\bibitem{manuscript} EGT/ECRAM manuscript draft, \texttt{project\_ecram/ref/Draft.pdf}
(Results \& Discussion).
\bibitem{bazant2004} M.~Z.~Bazant, K.~Thornton, A.~Ajdari,
\textit{Diffuse-charge dynamics in electrochemical systems}, Phys.\ Rev.\ E
\textbf{70}, 021506 (2004).
\bibitem{mei2018} B.-A.~Mei \textit{et al.},
\textit{Physical Interpretations of Nyquist Plots for EDLC Electrodes and
Devices}, J.\ Phys.\ Chem.\ C \textbf{122}, 194--206 (2018).
\bibitem{scharfetter} D.~L.~Scharfetter, H.~K.~Gummel,
\textit{Large-signal analysis of a silicon Read diode oscillator}, IEEE Trans.\
Electron Devices \textbf{16}, 64--77 (1969).
\bibitem{bard} A.~J.~Bard, L.~R.~Faulkner, \textit{Electrochemical Methods:
Fundamentals and Applications}, 2nd ed., Wiley (2001).
\end{thebibliography}

\end{document}
